\chapter{Resultados y discusión}
\label{chap:resultados}

\ph{Mientras el Capítulo~\ref{chap:pruebas} verificaba que el sistema está
bien construido, este capítulo evalúa si resuelve el problema planteado en el
Capítulo~\ref{chap:intro} y en qué medida: presenta los resultados obtenidos,
los contrasta con los criterios de éxito fijados, los compara con los
trabajos previos y los interpreta críticamente. Es el capítulo donde el
tribunal busca la evidencia de lo conseguido: cada afirmación debe apoyarse
en una tabla o una figura.}

\section{Diseño de la evaluación}
\label{sec:diseno_evaluacion}

\ph{Antes de mostrar números, define qué se mide y cómo. Retoma la
Tabla~\ref{tab:criterios-exito} del Capítulo~\ref{chap:intro}: sus criterios
de éxito son el guion de esta evaluación, y cada uno debe quedar respondido
en la Sección~\ref{sec:resultados_objetivos}. Describe:}

\begin{itemize}
  \item \textbf{\ph{Entorno de ejecución}}: \ph{hardware (CPU, RAM\dots) y
        software (sistema operativo, versiones relevantes) sobre los que se
        han obtenido las medidas, para que sean reproducibles.}
  \item \textbf{\ph{Datos o casos de estudio}}: \ph{origen, tamaño y
        representatividad de los datos, escenarios o usuarios empleados en la
        evaluación.}
  \item \textbf{\ph{Repeticiones}}: \ph{si hay medidas de rendimiento, indica
        cuántas veces se repite cada medición y por qué: una sola ejecución
        está sujeta a variabilidad (carga del sistema, caché, red\dots) y no
        es concluyente.}
  \item \textbf{\ph{Agregación}}: \ph{cómo se resumen los datos: media y
        desviación típica, mediana, percentiles (p.\,ej.\ p95 para tiempos de
        respuesta)\dots y por qué ese estadístico es el adecuado. Toda media
        debe ir acompañada de su desviación típica o de un intervalo de
        confianza; y si se comparan alternativas, conviene apoyar la
        comparación en un test estadístico sencillo (t-test o Wilcoxon).}
\end{itemize}

\ph{Rigor: distingue siempre lo que se ha \emph{medido} de lo que es una
\emph{impresión subjetiva}. «La interfaz es intuitiva» no es un resultado;
«8 de 10 usuarios completaron la tarea sin ayuda en menos de 2 minutos» sí
lo es.}

La evaluación se centra en comprobar la capacidad del generador para producir
modelos UVL con diferentes niveles de expresividad, la influencia de los
parámetros de configuración sobre los modelos obtenidos, la incorporación de
construcciones avanzadas del lenguaje, la integración dentro del ecosistema
UVLHub y Flamapy, y la disponibilidad de una batería de pruebas automatizadas
que permita validar su funcionamiento.

Durante el proceso de desarrollo se han utilizado diferentes configuraciones del generador y se han analizado los modelos UVL resultantes,
comprobando que las características seleccionadas durante la configuración se
reflejan correctamente en los artefactos generados. Estas comprobaciones se han
complementado con la ejecución de la suite de pruebas automatizadas y con
validaciones manuales de los principales flujos de uso de la aplicación.

La evaluación se ha realizado sobre el siguiente entorno:

\begin{itemize}
    \item \textbf{Sistema operativo:} Ubuntu 24.04.4 LTS.
    \item \textbf{Hardware:} Slimbook Executive con procesador
    Intel\textsuperscript{\textregistered} Core\textsuperscript{\texttrademark}
    i7-13700H (20 núcleos) y 32 GB de memoria RAM.
    \item \textbf{Navegador:} Firefox 148.0.2 (64 bits).
    \item \textbf{Lenguaje de desarrollo:} Python 3.13.11.
    \item \textbf{Contenedores:} Docker 29.7.2.
\end{itemize}

Adicionalmente, se realizará una evaluación del rendimiento del generador
mediante la medición del tiempo de inicialización del entorno y del tiempo
empleado en la generación de modelos una vez cargadas las dependencias. Para
ello, se utilizará una configuración con un volumen elevado de modelos generados,
con el objetivo de obtener una estimación representativa del comportamiento del
sistema en escenarios de mayor carga.

\section{Resultados}
\label{sec:resultados_objetivos}

\ph{Organiza esta sección con una subsección por objetivo (o por criterio de
éxito), en el mismo orden en que se enunciaron en el
Capítulo~\ref{chap:intro}. En cada subsección sigue el patrón: (1) presenta
el resultado en una tabla o figura, (2) referénciala desde el texto y (3)
interprétala frente al umbral fijado. Recuerda el principio: toda tabla o
figura se referencia desde el texto y se interpreta — un resultado sin
interpretación no es un resultado. Si el criterio de un objetivo ya quedó
verificado en un capítulo anterior (columna Cap.\ de la
Tabla~\ref{tab:criterios-exito}), su subsección se limita a recapitular esa
evidencia con \textbackslash ref y a dar el veredicto frente al criterio,
sin repetir datos. Las dos subsecciones siguientes modelan ambos patrones:
O1 ejemplifica la recapitulación y O2 la medición completa.}

\subsection{Resultados sobre el objetivo O1}
\label{subsec:resultados_o1}

\ph{Patrón RECAPITULACIÓN, para objetivos cuyo criterio ya quedó verificado
en un capítulo anterior según la columna Cap.\ de la
Tabla~\ref{tab:criterios-exito}. Siguiendo el ejemplo conductor, el criterio
de O1 (cobertura de todos los RF de prioridad alta) quedó verificado en los
Capítulos~\ref{chap:requisitos} y~\ref{chap:pruebas}: recapitula esa
evidencia con \textbackslash ref —la matriz de trazabilidad de la
Tabla~\ref{tab:trazabilidad_rf_pruebas} muestra que cada RF tiene al menos
una prueba que lo cubre, y la Tabla~\ref{tab:resumen_suite} acredita que la
suite completa pasa— y cierra con el veredicto frente al criterio («el
criterio de éxito asociado al objetivo O1 se considera cumplido»), sin
repetir datos ni añadir tablas nuevas.}

El objetivo O1 consistía en desarrollar un generador automático capaz de
producir modelos UVL con diferentes niveles de complejidad y expresividad,
incluyendo construcciones avanzadas del lenguaje que permitan ampliar la
variedad de modelos disponibles para el análisis de variabilidad.

Para verificar este objetivo se han generado modelos utilizando distintas
configuraciones del generador, comprobando que los parámetros seleccionados se
reflejan correctamente en los modelos UVL obtenidos y que estos mantienen una
estructura válida conforme al lenguaje.

En las siguientes imágenes se puede observar que efectivamente se ha alcanzado el objetivo definido. En la primera se puede ver la configuración usada para generar los modelos con expresividad avanzada, donde se han activado todos los \textit{major levels} junto con todos sus \textit{minor levels}

\begin{figure}[H]
    \centering
    \includegraphics[width=0.40\textwidth]{configuracion_verificar_O1.png}
    \caption{Configuración usada para verificar O1. Elaboración propia}
    \label{fig:configuracion-O1}
\end{figure}

Las dos siguientes imágenes muestran algunos de los modelos en cuestión, donde se puede comprobar que existen características, atributos y constraints de varios niveles de expresividad.

\begin{figure}[H]
    \centering
    \includegraphics[width=0.85\textwidth]{ejemplo1_verificar_O1.png}
    \caption{Ejemplo 1 para verificar O1. Elaboración propia}
    \label{fig:ejemplo1-modelo-uvl}
\end{figure}

\begin{figure}[H]
    \centering
    \includegraphics[width=0.85\textwidth]{ejemplo2_verificar_O1.png}
    \caption{Ejemplo 2 para verificar O1. Elaboración propia}
    \label{fig:ejemplo2-modelo-uvl}
\end{figure}

Los resultados obtenidos permiten verificar el cumplimiento del criterio de
éxito asociado al objetivo O1.

\subsection{Resultados sobre el objetivo O2}
\label{subsec:resultados_o2}

El objetivo O2 consistía en desarrollar un sistema de configuración flexible que
permitiera modificar los parámetros principales del proceso de generación y
adaptar los modelos obtenidos a diferentes escenarios.

Para verificar este objetivo se ha analizado el flujo de configuración del
generador, comprobando que los distintos parámetros disponibles se encuentran
organizados en diferentes etapas según su funcionalidad. La estructura del
asistente de configuración y las diferentes vistas que permiten definir estos
parámetros se muestran en la Sección~\ref{subsec:mapa_navegacion}.

Además, se ha comprobado durante el desarrollo que las modificaciones realizadas
sobre estos parámetros son aplicadas correctamente durante el proceso de
generación, lo cual a su vez puede comprobarse en el ~\ref{subsec:resultados_o1}.

Los resultados obtenidos permiten verificar el cumplimiento del criterio de
éxito asociado al objetivo O2.

\subsubsection*{Resultados sobre el objetivo O3}
\label{subsec:resultados_o3}

El objetivo O3 consistía en integrar la solución desarrollada dentro del
ecosistema Flamapy y permitir su utilización desde UVLHub como herramienta de
generación de modelos de características.

La prueba de que se ha alcanzado este objetivo es la propia integración del generador como
un plugin de Flamapy, cuya estructura propia puede comprobarse observando la imagen correspondiente en la Sección \ref{sec:tecnologias} (\textbf{Plugin generador basado en Flamapy}), así como su correcta utilización desde el flujo proporcionado por UVLHub.

Además, se ha verificado que el módulo desarrollado mantiene la compatibilidad
con las dependencias y componentes existentes del ecosistema, permitiendo su
ejecución dentro del entorno habitual de la aplicación.

Los resultados obtenidos permiten verificar el cumplimiento del criterio de
éxito asociado al objetivo O3.

\subsubsection*{Resultados sobre el objetivo O4}
\label{subsec:resultados_o4}

El objetivo O4 consistía en evaluar el comportamiento temporal del generador
en escenarios de generación con un volumen elevado de modelos de características.

Para realizar esta evaluación se ha utilizado una configuración con 1000 modelos
a generar, incorporando diferentes niveles de expresividad del lenguaje UVL. La
configuración empleada no representa un caso extremo en cuanto al tamaño
individual de los modelos, ya que se han utilizado rangos moderados de
características y restricciones, permitiendo obtener una estimación realista del
coste de generación en un escenario habitual con un volumen elevado de modelos.

La configuración utilizada durante la evaluación se muestra en la Figura
\ref{fig:configuracion-tiempo-generacion}.

\begin{figure}[H]
      \centering
      \makebox[\textwidth][c]{%
            \includegraphics[width=1.1\textwidth]{resultados_prueba_O4.png}
      }
      \caption{Resultado temporal de la generación de modelos. Elaboración propia}
      \label{fig:configuracion-tiempo-generacion}
\end{figure}

Los tiempos obtenidos durante la ejecución se dividen en dos fases diferenciadas:
la carga inicial de las dependencias necesarias para ejecutar el generador y la
generación posterior del conjunto de modelos UVL.

\begin{table}[H]
\centering
\begin{tabular}{l c}
\toprule
\textbf{Fase} & \textbf{Tiempo obtenido} \\
\midrule
Carga inicial de dependencias e inicialización & 4.05 s \\
Generación de 1000 modelos UVL & 7.8 s \\
\bottomrule
\end{tabular}
\caption{Tiempos obtenidos durante la evaluación temporal del generador. Elaboración propia}
\label{tab:tiempos-generador}
\end{table}

Los resultados obtenidos muestran que el coste asociado a la generación de un
volumen elevado de modelos es reducido una vez finalizada la fase inicial de
carga e inicialización. En este escenario, la generación de 1000 modelos UVL se
completa en 8.4 segundos, mostrando un comportamiento adecuado para el uso
previsto del generador. No obstante, estos tiempos dependen de la complejidad de
los modelos solicitados, pudiendo incrementarse en configuraciones con un mayor
número de características, restricciones o atributos. Además, los valores
obtenidos pueden variar en función de los recursos computacionales disponibles
en el equipo donde se ejecute el generador.

\subsubsection*{Resultados sobre el objetivo O5}
\label{subsec:resultados_o5}

El objetivo O5 consistía en validar el correcto funcionamiento del generador
mediante una batería de pruebas automatizadas que cubriese tanto la lógica
interna del motor de generación como el flujo de configuración integrado en
UVLHub.

Para verificar este objetivo se ha ejecutado la suite de pruebas del módulo
\textit{generator} dentro del entorno Docker de la aplicación. Estas pruebas
validan diferentes componentes del sistema, incluyendo la generación de modelos
UVL a partir de distintas configuraciones, la aplicación de los parámetros del
asistente de generación, la persistencia del estado entre etapas del flujo, las
opciones de salida generadas y el comportamiento del generador ante diferentes niveles de
expresividad del lenguaje UVL, verificando que las configuraciones relativas a
atributos, restricciones, cardinalidades y operadores avanzados producen los
resultados esperados.

La ejecución completa de la suite mediante \textit{rosemary} finalizó
correctamente, obteniéndose 134 pruebas superadas. En la Figura
\ref{fig:tests-exitosos-O5} se muestra el resultado de dicha ejecución.

\begin{figure}[H]
    \centering
    \includegraphics[width=1.1\textwidth]{salida_tests_O5.png}
    \caption{Resultado de la ejecución de la suite de pruebas automatizadas. Elaboración propia}
    \label{fig:tests-exitosos-O5}
\end{figure}

Los resultados obtenidos permiten verificar el cumplimiento del criterio de
éxito asociado al objetivo O5, demostrando la estabilidad de las funcionalidades
desarrolladas mediante pruebas automatizadas.


\section{Comparación con trabajos previos}
\label{sec:comparacion_trabajos}

\ph{Retoma la Tabla~\ref{tab:comparativa_relacionados} del
Capítulo~\ref{chap:estado_arte}, ahora con tus resultados en la mano: ¿qué
celdas marcadas como «Parcial» o «No» en los otros trabajos cubre ya tu
solución, con qué evidencia de la Sección~\ref{sec:resultados_objetivos}? ¿Y
en qué queda tu trabajo por detrás de ellos? Una comparación honesta en ambas
direcciones es lo que distingue un análisis riguroso de un alegato.}

\ph{Si los trabajos relacionados publican métricas comparables (mismos datos
o condiciones equiparables), añade aquí una tabla numérica con una fila por
métrica y una columna por trabajo, citando la fuente de cada
cifra~\cite{apellido2024articulo}. Las cifras propias deben ir acompañadas
de su dispersión (desviación típica o intervalo de confianza), no solo de la
media. Si no existen métricas comparables, no la
fuerces: justifica por qué la comparación directa no es posible (datos,
entorno o alcance distintos) y mantén la comparación cualitativa por
criterios.}

La Tabla~\ref{tab:comparativa-trabajos} del Capítulo~\ref{chap:estado-arte}
recoge una comparación cualitativa entre diferentes herramientas relacionadas
con la generación automática de modelos de características. En ella se
consideran tres aspectos principales: la capacidad de generación automática, el
grado de personalización ofrecido y la incorporación de construcciones de
expresividad avanzada.

Los resultados obtenidos durante la evaluación del generador desarrollado
permiten observar que la solución propuesta cubre los tres criterios analizados.
A diferencia de S.P.L.O.T. y FeatureIDE, cuyas capacidades de generación están
principalmente orientadas a modelos basados en construcciones booleanas,
el generador desarrollado permite producir modelos UVL incorporando elementos
adicionales como atributos tipados, restricciones aritméticas, restricciones
sobre cadenas, funciones agregadas y diferentes mecanismos de cardinalidad.

En comparación con BeTTy, la solución desarrollada mantiene la capacidad de
configuración flexible de los parámetros de generación, pero amplía el rango de
expresividad de los modelos generados. Mientras que BeTTy permite definir
escenarios experimentales con diferentes niveles de complejidad estructural, el
generador desarrollado incorpora además construcciones propias de lenguajes de
variabilidad más recientes como UVL.

Además, la integración realizada dentro del ecosistema Flamapy y UVLHub supone
una diferencia respecto a las herramientas analizadas, ya que permite utilizar
el generador directamente dentro de una plataforma web orientada a la gestión y
análisis de modelos de características, manteniendo la compatibilidad con las
herramientas existentes del ecosistema.

No se ha realizado una comparación cuantitativa directa en términos de tiempo de
generación o volumen máximo de modelos, debido a que las herramientas estudiadas
no proporcionan escenarios de evaluación homogéneos que permitan establecer una
comparación justa, por lo que la comparación se ha centrado en las
capacidades funcionales y de expresividad.

En consecuencia, los resultados obtenidos muestran que el trabajo desarrollado
amplía las capacidades de los generadores existentes, especialmente en cuanto a
la generación de modelos con mayor expresividad y su integración dentro del
ecosistema UVLHub, lo cual permite afirmar que efectivamente el trabajo realizado en este TFG no queda por detrás del resto a pesar de que herramientas como BeTTy cuenten con una mayor madurez histórica y una validación más extensa en escenarios experimentales.

\section{Discusión}
\label{sec:discusion}

\ph{Interpretación global, sin repetir los números anteriores: (1) ¿qué
significan los resultados en conjunto? ¿responden al problema de la
Sección~\ref{sec:problema}? (2) ¿qué resultados fueron inesperados y a qué
los atribuyes (hipótesis razonadas, no excusas)? (3) ¿qué compromisos
(\emph{trade-offs}) revelan: rendimiento frente a simplicidad, coste frente a
precisión, generalidad frente a especialización\dots?}

\ph{Cierra con un párrafo breve de amenazas a la validez: qué podría
debilitar estas conclusiones, distinguiendo validez \emph{interna} (factores
del propio experimento: datos sintéticos, selección de casos, configuración),
\emph{externa} (hasta qué punto generalizan los resultados a otros entornos,
datos o usuarios) y \emph{de constructo} (¿miden las métricas elegidas
realmente lo que se quiere evaluar?). Basta un párrafo: el tratamiento
completo de las limitaciones se hace en la Sección~\ref{sec:limitaciones},
no lo dupliques aquí.}

Los resultados obtenidos muestran que la solución desarrollada responde al
problema planteado inicialmente: la necesidad de disponer de un generador capaz
de producir modelos de características con una mayor expresividad que las
aproximaciones tradicionales basadas principalmente en construcciones
booleanas. La integración del generador dentro del ecosistema Flamapy y UVLHub
permite disponer de una herramienta flexible para la generación de modelos UVL
destinados a tareas posteriores de análisis y experimentación.

Durante el desarrollo se encontraron diferentes dificultades relacionadas con
la integración y adaptación de los componentes existentes. Sin embargo, la
evaluación final confirmó que el generador soporta correctamente las
funcionalidades planteadas inicialmente, incluyendo los diferentes niveles de
expresividad y las distintas configuraciones disponibles.

La incorporación de construcciones avanzadas implica ciertos compromisos. Por
un lado, aumentar la expresividad mediante atributos tipados, restricciones
aritméticas, funciones agregadas o restricciones sobre cadenas incrementa la
complejidad interna del generador. Por otro lado, ofrecer una configuración más
flexible permite adaptar la generación a distintos escenarios, aunque aumenta
el número de parámetros que el usuario debe comprender. Asimismo, una mayor
riqueza en los modelos generados puede suponer un incremento del coste
computacional en configuraciones más complejas.

Finalmente, existen algunas amenazas a la validez de los resultados. La
evaluación se ha realizado utilizando modelos sintéticos y configuraciones
diseñadas para cubrir las funcionalidades principales del generador, por lo que
los resultados podrían variar en otros escenarios. Además, los tiempos
obtenidos dependen del entorno de ejecución y no se ha realizado una
comparación cuantitativa directa con otras herramientas bajo condiciones
equivalentes.

% --- Transición al siguiente capítulo ---
\ph{Cierra el capítulo con un párrafo puente de 2--4 líneas: qué queda
establecido aquí y qué pregunta abre el capítulo siguiente. Por ejemplo:
«Evaluados los resultados frente a los criterios de éxito y discutido su
alcance, el Capítulo~\ref{chap:conclusiones} sintetiza el grado de
cumplimiento de los objetivos planteados en el Capítulo~\ref{chap:intro},
las limitaciones del trabajo y las líneas de trabajo futuro que estos
resultados abren».}

Los resultados obtenidos y la discusión realizada permiten establecer el grado de
cumplimiento de los objetivos planteados y validar la solución desarrollada
frente a los criterios de éxito definidos. A partir de este análisis, el
Capítulo~9 aborda las conclusiones finales del trabajo, las limitaciones
identificadas y las posibles líneas de trabajo futuro derivadas de los
resultados obtenidos.